% Options for packages loaded elsewhere
\PassOptionsToPackage{unicode}{hyperref}
\PassOptionsToPackage{hyphens}{url}
%
\documentclass[
  11pt,
]{article}
\usepackage{amsmath,amssymb}
\usepackage{iftex}
\ifPDFTeX
  \usepackage[T1]{fontenc}
  \usepackage[utf8]{inputenc}
  \usepackage{textcomp} % provide euro and other symbols
\else % if luatex or xetex
  \usepackage{unicode-math} % this also loads fontspec
  \defaultfontfeatures{Scale=MatchLowercase}
  \defaultfontfeatures[\rmfamily]{Ligatures=TeX,Scale=1}
\fi
\usepackage{lmodern}
\ifPDFTeX\else
  % xetex/luatex font selection
\fi
% Use upquote if available, for straight quotes in verbatim environments
\IfFileExists{upquote.sty}{\usepackage{upquote}}{}
\IfFileExists{microtype.sty}{% use microtype if available
  \usepackage[]{microtype}
  \UseMicrotypeSet[protrusion]{basicmath} % disable protrusion for tt fonts
}{}
\makeatletter
\@ifundefined{KOMAClassName}{% if non-KOMA class
  \IfFileExists{parskip.sty}{%
    \usepackage{parskip}
  }{% else
    \setlength{\parindent}{0pt}
    \setlength{\parskip}{6pt plus 2pt minus 1pt}}
}{% if KOMA class
  \KOMAoptions{parskip=half}}
\makeatother
\usepackage{xcolor}
\usepackage[margin=1in]{geometry}
\setlength{\emergencystretch}{3em} % prevent overfull lines
\providecommand{\tightlist}{%
  \setlength{\itemsep}{0pt}\setlength{\parskip}{0pt}}
\setcounter{secnumdepth}{-\maxdimen} % remove section numbering
\ifLuaTeX
  \usepackage{selnolig}  % disable illegal ligatures
\fi
\usepackage{bookmark}
\IfFileExists{xurl.sty}{\usepackage{xurl}}{} % add URL line breaks if available
\urlstyle{same}
\hypersetup{
  hidelinks,
  pdfcreator={LaTeX via pandoc}}

\author{}
\date{}

\begin{document}

\section{Exact edge general position numbers of higher-order Fibonacci
cubes}\label{exact-edge-general-position-numbers-of-higher-order-fibonacci-cubes}

Ryutaro Yonezu\\
Independent Researcher\\
Version 1.0 --- 2026-11-19

\textbf{Preprint. NOT HUMAN PEER REVIEWED.} Independent AI mathematics
audit: \textbf{THEOREM CONDITIONAL PASS / UNIQUENESS CONDITIONAL PASS}.
NOVELTY INCONCLUSIVE / REFERENCES PARTIAL. All rights reserved.

This is an analytic proof. Finite computation provides corroboration
only. The theorem and complete classification of maximum sets below
reproduce the frozen v0.1 mathematical manuscript without changes. The
conditional AI audit status is based on an author-provided transcription
of an independent AI response. An independent verbatim comparison with
the source chat was not performed; the auditor's independent checking
code was not obtained.

\subsection{Model and theorem}\label{model-and-theorem}

For integers k\textgreater=3 and n\textgreater=2 let G=Q\_n(1\^{}k) be
the subgraph of the binary n-cube induced by the words without k
consecutive ones. Write A\_n=A\_n\^{}(k) for their number, A\_0=1. An
edge general position set S has no three distinct edges on a common
shortest path. Denote its maximum cardinality by gp\_e(G). Let Theta\_i
be the edges flipping coordinate i.

\textbf{Theorem.} For all k\textgreater=3 and n\textgreater=2,
gp\_e(G)=2(A\_n-A\_(n-1)). If n\textgreater=k+1 the unique maximum set
is Theta\_1 union Theta\_n. If 2\textless=n\textless=k, precisely the
unions of two distinct coordinate classes are maximum sets.

No statement for k=2 is assumed: the proof is independent of No.75 and
No.76.

\subsection{1. Isometry and coordinate
counts}\label{1-isometry-and-coordinate-counts}

The legal-word family is downward closed in the coordinatewise order.
Between any legal x,y, go from x down to x AND y, then up to y, flipping
each differing bit once. Every intermediate word is legal, hence graph
distance equals Hamming distance. In particular a geodesic never repeats
a coordinate; a union of two coordinate classes is in edge general
position.

Every legal word with bit i=1 corresponds bijectively to its downward
edge in Theta\_i. The words with bit i=0 split independently into two
legal words, so

\textbar Theta\_i\textbar{} = A\_n - A\_(i-1) A\_(n-i).

Concatenation embeds the legal words of length n-1 into the pairs of
legal words of lengths i-1,n-i. Consequently A\_(i-1)
A\_(n-i)\textgreater=A\_(n-1). For i=1,n equality holds, and
B=2(A\_n-A\_(n-1)) is a lower bound for gp\_e(G).

If n\textgreater=k+1 and 1\textless i\textless n, choose positive
integers a\textless=i-1,b\textless=n-i with a,b\textless=k-1 and a+b=k.
Such integers exist since both lengths are positive and their sum is
n-1\textgreater=k. A legal left word ending in a ones and a legal right
word starting in b ones (all other bits zero) have illegal
concatenation. Thus the inequality is strict for every interior i. When
n\textless=k, A\_(n-1)=2\^{}(n-1) and both factors are powers of two, so
all coordinate classes have equal size.

\subsection{2. Two-arm projection
lemma}\label{2-two-arm-projection-lemma}

In any isometric induced subgraph of a hypercube, suppose a selected
star at c contains edges in distinct coordinates i,j. A selected edge e
in a coordinate l different from i,j cannot have projection on i,j equal
to either of the two leaves c\^{}i or c\^{}j (superscript means flip
that coordinate).

Proof: choose that leaf z, and the endpoint w of e with w\_l=z\_l. By
hypothesis z and w agree on i,j,l, so any shortest path joining z and w
flips none of these coordinates. Concatenate the other leaf
-\textgreater{} c -\textgreater{} z, a shortest z-w path, and e. All
coordinate flips occur at most once, and because graph distances are
Hamming distances this is a geodesic containing three selected edges.
Contradiction.

\subsection{3. Star forest lemma}\label{3-star-forest-lemma}

If S uses at least three distinct coordinates, the graph whose edges are
S is a star forest (ignore isolated vertices).

Proof: if a connected component contains a simple three-edge path, its
directions must repeat (otherwise the path is a geodesic). As
consecutive directions cannot repeat in a simple path, the pattern is
i,j,i. The two consecutive two-arm stars have respective leaf
projections covering all four binary patterns on i,j (the four vertices
form three sides of an ambient square). A selected edge in a third
coordinate is therefore forbidden by Section 2. This contradicts the
hypothesis. Hence there is no simple path of length three. The graph is
bipartite. Any connected bipartite graph with no simple three-edge path
is a star: a cycle would contain such a path, and a tree of diameter at
most two is a star.

In particular, with largest star degree d,
\textbar S\textbar\textless=d/(d+1) A\_n.

\subsection{4. Uniform density of the two end
classes}\label{4-uniform-density-of-the-two-end-classes}

For n\textless k, A\_n=2\^{}n and B=A\_n. For n=k, A\_n=2\^{}k-1 and
B=2\^{}k-2. For n\textgreater=k+1, counting legal words of length n-1
extended by one bit gives A\_n=2 A\_(n-1)-A\_(n-k-1): the forbidden
extensions of otherwise legal words are exactly the words ending in 0
followed by k ones, with an arbitrary legal prefix of length n-k-1.
Hence B=A\_n-A\_(n-k-1).

Claim: B\textgreater=10 A\_n/11 for all n\textgreater=4,
k\textgreater=3, with strict inequality for n\textgreater=10.

For k\textgreater=4 and n\textgreater=k+1, from each legal prefix of
length n-k-1 append a zero and any legal word of length k. This gives
A\_n \textgreater= (2\^{}k-1) A\_(n-k-1) \textgreater=15 A\_(n-k-1). The
cases n\textless=k are immediate from the explicit counts
(n=k\textgreater=4). For k=3, A\_n satisfies
A\_n=A\_(n-1)+A\_(n-2)+A\_(n-3) for n\textgreater=3. The inequality
A\_n\textgreater=11 A\_(n-4) at n=4,5,6 follows from
A\_4=13,A\_5=24,A\_6=44 and A\_0=1,A\_1=2,A\_2=4. The difference
satisfies the same recurrence from n=7 onward, hence is strictly
positive for every n\textgreater=7. This proves the claim.

\subsection{5. Large stars exclude a fixed fraction of
vertices}\label{5-large-stars-exclude-a-fixed-fraction-of-vertices}

Suppose d\textgreater=11. Among the d coordinates of a largest star at
c, choose six pairwise nonconsecutive coordinates: sorting the d
coordinates and selecting alternating ones gives at least
ceil(d/2)\textgreater=6. Split the chosen coordinates into two triples.

Let X be a uniformly random legal vertex. Conditional on any specified
values of some other chosen coordinates (with nonempty conditioning
event), for any still-unfixed chosen coordinate i we have

1/5 \textless= Pr(X\_i=1 \textbar{} condition) \textless= 1/2.

Indeed, deleting a one at i is an injection from the conditioned words
with X\_i=1 into those with X\_i=0. Conversely, set bit i to one and
clear both immediate neighbors (only existing neighbors at boundaries).
This creates an isolated one, preserves legality for k\textgreater=2,
and preserves the conditioning since no chosen coordinate is adjacent to
i. This map from the zero class into the one class has at most four
preimages per image (the former neighboring bits). Thus \#ones
\textless= \#zeros \textless=4 \#ones. Both classes are nonempty. It
follows that the disagreement indicator Y\_i=1\_\{X\_i != c\_i\} has
conditional probabilities for either value at least 1/5 and at most 4/5.

In either triple, conditional on the first disagreement indicator and
any earlier specified chosen bits, the probability that both remaining
indicators equal the first is at most (4/5)\^{}2. Thus the probability
that the triple is mixed (contains both zero and one) is at least 9/25.
Conditioning on every possible full assignment of the first triple, the
same bound holds for the second. Therefore the probability that both
triples are mixed is at least (9/25)\^{}2=81/625.

For any vertex x in this event, its projection on the d star coordinates
has at least two agreements and at least two disagreements with c. For
any incident selected edge in direction l, remove l if it is a star
coordinate. At least one agreeing and one disagreeing star coordinate
remain; call them i,j. The edge has the projected pattern of one of the
two leaves c\^{}i,c\^{}j on this pair. Section 2 forbids that edge.
Hence none of these vertices is incident to S.

At least 81/625 of the vertices are excluded. As S is a star forest, its
number of edges is at most its number of incident vertices. In
particular \textbar S\textbar{} \textless= 544 A\_n/625 \textless{} 10
A\_n/11 \textless= B (n\textgreater=11).

\subsection{6. Completion and equality}\label{6-completion-and-equality}

If S uses at most two coordinates, Section 1 bounds its size by B;
equality requires two complete largest coordinate classes. If S uses at
least three, it is a star forest. When d\textgreater=11 Section 5 gives
a strict bound below B. When d\textless=10,
\textbar S\textbar\textless=d A\_n/(d+1)\textless=10 A\_n/11. For
n\textgreater=4 the bound is strict below B: if n\textless=9 then
d\textless=n\textless=9 and the degree ratio is strictly below 10/11; if
n\textgreater=10 Section 4 gives B\textgreater10 A\_n/11. For n=3,
d\textless=3 and B is 6 for k=3, 8 for k\textgreater3; 3 A\_n/4
\textless{} B in both cases. For n=2 at least three coordinates are
impossible. These cases exhaust all sets. The equality classification
follows from Section 1. QED.

\subsection{Reproduction and
limitations}\label{reproduction-and-limitations}

The proof is analytic and does not depend on finite computational cases.
verify.py uses only the Python standard library (Python 3.10+) and
independently enumerates legal words, conditional counts, projection
events, and small forbidden-triple optimization problems. These checks
corroborate the proof; none substitutes for an independent review of its
universal arguments.

\subsection{Prior work (preliminary)}\label{prior-work-preliminary}

S. Klavzar and E. Tan, Edge General Position Sets in Fibonacci and Lucas
Cubes, Bull. Malays. Math. Sci. Soc. 46 (2023), 120, DOI
10.1007/s40840-023-01517-y. Known context for edge general position in
partial cubes and ordinary Fibonacci and Lucas cubes. Two-coordinate
lower bound is a standard partial-cube observation, not claimed new.

The size of k-th order generalized Fibonacci cubes, arXiv:2601.17273
(2026). Relevant counting context. Direct formula overlap not located in
preliminary title/abstract search; full-text review pending. No priority
claim.

\end{document}
